\HMTNarrativeHeading{La unidad conservada en sus realizaciones}

El recorrido del artículo vuelve a su tesis inicial con operaciones
ya desarrolladas. APP conserva la evaluación y su cociente;
TRIT distingue orientación y régimen; el TPK compone transporte,
selección y prolongación. La conexión nonádica actualiza la
historia durante el retorno de fase. El refinamiento produce
las lecturas numéricas y conserva las incidencias que sostienen
su realización excepcional.

La constante holográfica de acoplamiento aritmético-geométrico
reúne tres funciones precisas. Su registro admite una codificación
posicional reversible; interviene en la clausura de alfa y en
la incidencia marcada; su órbita permite identificar operadores
con cotas de recuperación. La ley bilateral transporta esa
capacidad mediante la conservación del estado y de sus formas.

Las vacancias explican cómo la historia avanza cuando una
capacidad permanece fija. La sección electrónica muestra una
persistencia con trayectoria y registro. Las secciones de acción,
los canales orientados, los lectores de área y la transducción
térmica convierten los resultados anteriores en relaciones
físicas calculables. Las realizaciones excepcionales conservan
sus portadores, incidencias y mapas dimensionales.

La conservación evolutiva de la unidad dimensional queda
expresada en varios niveles enlazados: partición normalizada,
refinamiento reversible, transporte de formas, archivo con
terminal, recuperación estable y transformación de observables.
Cada nivel tiene su objeto y su prueba; las aplicaciones entre
ellos construyen la continuidad del argumento. La doble
proyección mantiene el centro de esa organización: valor e
incidencia proceden del mismo estado, cuya transformación
conserva relaciones recuperables.

